\section{ML-Based Bitstream Inspection}
\label{sec:ml-inspection}

We use convolutional neural networks (CNNs) to classify the fixed-size
bitstream representations introduced in Section~\ref{sec:method-dataset}. Prior work
has shown that bitstream-derived representations contain learnable structure and
can be used to distinguish benign and malicious FPGA designs
\cite{rettkowski2019inspection,chaudhuri2022detection,
elnaggar2023learning,palumbo2022trojan}. CNNs are also well suited to our
deployment goal because quantized models can be translated into FPGA circuits
using tools such as hls4ml~\cite{duarte2018fast,aarrestad2021fast}.

\subsection{CNN Architecture}
\label{sec:cnn-architecture}

The selected checker uses a seven-stage convolutional architecture operating on
a $256\times256$ byte-image tensor. Each stage consists of a quantized
convolution, quantized ReLU activation, and max pooling. Channel width increases
from 8 filters in the first stage to 32 filters in the later stages. A final
average-pooling operation reduces the remaining feature map before a single
quantized dense output produces the classification logit.

Table~\ref{tab:cnn-architecture} summarizes the architecture. We use 8-bit
weights and activations (W8A8) and do not use batch normalization. The output
logit is converted to a maliciousness probability using a sigmoid; the decision
threshold is applied later by the deployment logic.

\begin{table}[t]
\centering
\caption{Seven-stage quantized CNN architecture used by the on-FPGA
checker. All convolutional stages use quantized ReLU activations and W8A8
quantization.}
\label{tab:cnn-architecture}
\small
\begin{tabular}{lcccc}
\toprule
Stage & Filters & Kernel & Stride & Pooling \tabularnewline
\midrule
Conv 0 & 8 & $5\times5$ & $2\times2$ & $2\times2$ \tabularnewline
Conv 1 & 16 & $3\times3$ & $1\times1$ & $2\times2$ \tabularnewline
Conv 2 & 24 & $3\times3$ & $1\times1$ & $2\times2$ \tabularnewline
Conv 3 & 24 & $3\times3$ & $1\times1$ & $2\times2$ \tabularnewline
Conv 4 & 32 & $3\times3$ & $1\times1$ & $2\times2$ \tabularnewline
Conv 5 & 32 & $3\times3$ & $1\times1$ & $2\times2$ \tabularnewline
Conv 6 & 32 & $3\times3$ & $1\times1$ & $2\times2$ \tabularnewline
Output & 1 & -- & -- & AvgPool \tabularnewline
\bottomrule
\end{tabular}
\end{table}

\subsection{Training}
\label{sec:cnn-training}

Models are implemented in TensorFlow~\cite{abadi2016tensorflow} using
QKeras~\cite{coelho2020automatic}. We use quantization-aware training
\cite{coelho2020automatic}, so that the model is trained under the same
fixed-precision behavior used during hardware generation. Training minimizes
binary cross-entropy from logits. We augment the training data with horizontal
and vertical flips and use a balanced production training set containing equal
numbers of benign and malicious samples.

During model development, we use stratified cross-validation to compare model
configurations~\cite{kohavi1995study}. After selecting the final architecture, we train a production
model separately and evaluate it on the held-out test set described in
Section~\ref{sec:eval-setup}.

\subsection{Hardware-Aware Design-Space Exploration}
\label{sec:cnn-design-space}

Model selection considers hardware feasibility in addition to software
classification quality. Rather than exhaustively evaluating every combination
of parameters, we explore the design space in stages. We first vary input
resolution and CNN depth to determine how model capacity and retained bitstream
information affect classification quality. We then evaluate quantization and
pruning as model-compression techniques before passing selected networks to the
hardware synthesis flow described in Section~\ref{sec:on-fpga-checker}.

The principal model-level dimensions are therefore input resolution, network
depth, and numerical precision. Hardware parallelism is explored separately
after conversion to hls4ml by varying operation reuse and per-layer
configuration.

The resulting design-space exploration selects the $256\times256$, seven-layer
W8A8 CNN as the on-FPGA checker used throughout the remainder of the paper.
We also synthesize the same architecture at $512\times512$ to evaluate how
hardware cost scales with input resolution.